\section{风险与研究议程：开放训练下一步该做什么}

\subsection{当前最关键的三类研究缺口}
\begin{table}[H]
\centering
\renewcommand{\arraystretch}{1.2}
\begin{tabular}{p{3.0cm}p{4.0cm}p{4.2cm}}
\toprule
\textbf{缺口} & \textbf{现状问题} & \textbf{下一步方向} \\
\midrule
验证器生态 & 可验证任务覆盖仍窄，跨领域迁移弱 & 构建可组合 verifier 库与任务分层标准 \\
数据治理 & 合成数据质量波动，污染难以统一审计 & 统一去污染协议与开源数据审计工具链 \\
成本可控性 & 测试时扩展收益与成本预测不稳定 & 建立 ``质量-延迟-成本'' 联合调度策略 \\
\bottomrule
\end{tabular}
\caption{讲座内容可落地的三类研究议程}
\end{table}

\begin{knowledgebox}{对工程团队的直接启发}
将 ``训练配方'' 文档化与版本化，与代码同等管理。每次能力变化都应可追溯到数据、目标函数、优化器、推理策略中的具体改动。
\end{knowledgebox}

\subsection{给研究者和产品团队的两条建议}
\begin{importantbox}{建议 1：优先建设可复现实验底座}
即便短期不追最高分，也要先把训练、评估、回归测试跑通。缺乏可复现基础时，任何改进都可能不可持续。
\end{importantbox}

\begin{importantbox}{建议 2：把推理预算当作产品参数}
推理时扩展不是算法附属，而是产品策略。要按请求价值分层分配预算，而不是对所有请求使用同一解码策略。
\end{importantbox}

\begin{warningbox}{开放并不意味着低标准}
开放训练路线同样需要严格的数据合规、评估去污染和安全审查。否则开放会变成复制噪声，而不是复制进步。
\end{warningbox}

\subsection{本章小结}
下一阶段开放训练的胜负手，不是某个单一新算法，而是验证器生态、数据治理和推理预算调度这三件基础设施。


